\documentclass[11pt,oneside]{article}
% Standalone ProveIt research note: no external inputs are required.
\usepackage[a4paper,margin=26mm,headheight=15pt]{geometry}
\usepackage[T1]{fontenc}
\usepackage[utf8]{inputenc}
\IfFileExists{libertinus.sty}{\usepackage{libertinus}}{\usepackage{lmodern}}
\usepackage{microtype}
\usepackage{amsmath,amssymb,amsthm,mathtools,mathrsfs,bm}
\usepackage{booktabs,longtable,tabularx,array}
\usepackage{graphicx}
\usepackage{xcolor}
\usepackage{enumitem}
\usepackage{fancyhdr}
\usepackage{xurl}
\usepackage{listings}
\usepackage[most]{tcolorbox}
\usepackage{hyperref}
\usepackage[nameinlink,noabbrev,capitalise]{cleveref}

\definecolor{linkblue}{RGB}{24,72,124}
\definecolor{deepgreen}{RGB}{23,102,69}
\definecolor{deepred}{RGB}{140,42,42}
\definecolor{deepgold}{RGB}{122,91,19}
\definecolor{shadegray}{RGB}{246,247,249}
\definecolor{rulegray}{RGB}{195,201,208}
\definecolor{palegreen}{RGB}{233,245,243}
\definecolor{paleblue}{RGB}{236,243,251}
\definecolor{palegold}{RGB}{251,245,230}
\definecolor{palered}{RGB}{251,236,236}

\hypersetup{
  hypertexnames=false,
  colorlinks=true,
  linkcolor=linkblue,
  citecolor=deepgreen,
  urlcolor=linkblue,
  pdftitle={Complete Asymptotic Transseries for the Double Factorial and Its Inverses},
  pdfauthor={Research note prepared for Vladimir Reshetnikov},
  pdfsubject={Double factorial, Bernoulli and Bell coefficients, Lambert W inversion, parity branches, periodic interpolation},
  pdfkeywords={double factorial, asymptotic expansion, transseries, Lambert W, inverse function, Bernoulli polynomial, Bell polynomial, Lagrange inversion}
}

\pagestyle{fancy}
\fancyhf{}
\fancyhead[L]{\small Double-factorial transseries}
\fancyhead[R]{\small\nouppercase{\leftmark}}
\fancyfoot[C]{\thepage}
\renewcommand{\headrulewidth}{0.4pt}

\setcounter{tocdepth}{2}
\setcounter{secnumdepth}{3}
\setlist[itemize]{leftmargin=1.5em,itemsep=0.3ex,topsep=0.7ex}
\setlist[enumerate]{leftmargin=1.9em,itemsep=0.4ex,topsep=0.7ex}
\allowdisplaybreaks[2]
\emergencystretch=2em
\raggedbottom

\theoremstyle{plain}
\newtheorem{theorem}{Theorem}[section]
\newtheorem{proposition}[theorem]{Proposition}
\newtheorem{lemma}[theorem]{Lemma}
\newtheorem{corollary}[theorem]{Corollary}
\theoremstyle{definition}
\newtheorem{definition}[theorem]{Definition}
\newtheorem{algorithm}[theorem]{Algorithm}
\newtheorem{example}[theorem]{Example}
\theoremstyle{remark}
\newtheorem{remark}[theorem]{Remark}

\crefname{algorithm}{algorithm}{algorithms}
\Crefname{algorithm}{Algorithm}{Algorithms}

\newtcolorbox{mainbox}[1][]{
  enhanced,breakable,colback=paleblue,colframe=linkblue,
  boxrule=0.7pt,arc=1.2mm,left=2mm,right=2mm,top=1.4mm,bottom=1.4mm,
  title={Main result},fonttitle=\bfseries,#1}
\newtcolorbox{warningbox}[1][]{
  enhanced,breakable,colback=palegold,colframe=deepgold,
  boxrule=0.7pt,arc=1.2mm,left=2mm,right=2mm,top=1.4mm,bottom=1.4mm,
  title={Important distinction},fonttitle=\bfseries,#1}
\newtcolorbox{summarybox}[1][]{
  enhanced,breakable,colback=palegreen,colframe=deepgreen,
  boxrule=0.7pt,arc=1.2mm,left=2mm,right=2mm,top=1.4mm,bottom=1.4mm,
  title={Summary},fonttitle=\bfseries,#1}
\newtcolorbox{checkbox}[1][]{
  enhanced,breakable,colback=shadegray,colframe=rulegray,
  boxrule=0.6pt,arc=1.0mm,left=2mm,right=2mm,top=1.2mm,bottom=1.2mm,
  title={Check},fonttitle=\bfseries,#1}

\lstdefinestyle{code}{
  basicstyle=\ttfamily\small,
  backgroundcolor=\color{shadegray},
  frame=single,
  rulecolor=\color{rulegray},
  breaklines=true,
  columns=fullflexible,
  keepspaces=true,
  showstringspaces=false,
  xleftmargin=0.4em,
  xrightmargin=0.4em,
  aboveskip=0.8em,
  belowskip=0.8em
}
\lstset{style=code}

\newcommand{\e}{\mathrm e}
\newcommand{\ii}{\mathrm i}
\newcommand{\dd}{\,\mathrm d}
\newcommand{\W}{W}
\newcommand{\R}{\mathbb R}
\newcommand{\Q}{\mathbb Q}
\newcommand{\N}{\mathbb N}
\newcommand{\BigO}{\mathcal O}
\newcommand{\coeff}[2]{[\,#1^{#2}\,]}
\newcommand{\BellP}{\mathrm B}
\newcommand{\BellC}{Y}
\newcommand{\stirlingone}[2]{\genfrac{[}{]}{0pt}{}{#1}{#2}}
\newcommand{\DF}{\mathcal D}
\newcommand{\Dc}{\mathcal D_{\!c}}
\newcommand{\sgn}{\operatorname{sgn}}

\title{\textbf{Complete Asymptotic Transseries for the Double Factorial and Its Inverses}\\[0.4em]
\large Bernoulli--Bell coefficient formulae, Lambert-
\(W\) reversion, parity branches, and the periodic OEIS interpolation}
\author{Research note prepared for Vladimir Reshetnikov}
\date{3 September 2026}

\begin{document}
\maketitle

\begin{abstract}
The double factorial is an interleaving of two gamma-type sequences.  That
apparently elementary parity split is decisive for inversion: there is no
canonical real inverse of a discrete sequence until an interpolation is fixed.
This article first treats the two exact parity branches
\(
 \DF_\delta(x)=(2/\pi)^{\delta/2}2^{x/2}\Gamma(x/2+1)
\), \(\delta\in\{0,1\}\), and then the smooth cosine interpolation recorded in
OEIS~A006882.  With the optimal centre \(X=x+1\), both parity branches have the
same logarithmic Stirling tail.  Its coefficients are
\[
 c_k=\frac{(1-2^{2k-1})B_{2k}}{2k(2k-1)},
\]
and all multiplicative coefficients are complete exponential Bell
polynomials.  For the large inverse, the dominant equation is solved exactly
by a Lambert-
\(W\) core.  The remaining displacement contains only
\(T^{-1},T^{-3},\ldots\); its coefficient functions are rational polynomials in
\(\Lambda^{-1}\), generated to arbitrary order by a triangular coefficient
recurrence or, explicitly, by partial Bell polynomials.  The first five blocks
are displayed and a scaling identity connects them to the half-shifted inverse
Gamma coefficients.  Expanding the Lambert core gives the full outer
iterated-logarithmic layer with general unsigned-Stirling-number coefficients.
For the cosine interpolation, a bounded periodic factor creates an additional
oscillatory \(1/\Lambda\)-sector; Lagrange--B\"urmann inversion gives its carrier
in closed Fourier form, and a second triangular recurrence gives all algebraic
corrections.  Fixed-order remainders and numerical residual checks are
included throughout.
\end{abstract}

\tableofcontents

\section{Introduction and scope}
\label{sec:introduction}

For integers \(n\geq0\), the double factorial sequence is determined by
\begin{equation}
 a_n=n!!,\qquad a_0=a_1=1,\qquad a_n=n\,a_{n-2}.
 \label{eq:sequence-recurrence}
\end{equation}
It is OEIS~A006882 \cite{OEIS}.  The leading estimate
\[
 n!!\sim
 \begin{cases}
 \sqrt\pi\,n^{(n+1)/2}\e^{-n/2},&n\ \hbox{even},\\
 \sqrt2\,n^{(n+1)/2}\e^{-n/2},&n\ \hbox{odd},
 \end{cases}
\]
is classical, but it conceals several structures needed for a genuinely
all-orders inverse:
\begin{itemize}
\item the centre \(X=n+1\), rather than \(n\), removes all even powers from the
      logarithmic correction;
\item parity affects only one normalization constant after this centring;
\item inversion of \(X(\log X-1)/2\) is naturally expressed by \(W_0\);
\item the inverse Stirling corrections live on the sparse grid
      \(T^{-1},T^{-3},T^{-5},\ldots\);
\item a smooth interpolation that merges both parities necessarily introduces
      extra information, and its inverse need not have the same transseries as
      either parity branch.
\end{itemize}

The coefficient calculus used below follows the Bell-polynomial and formal
reversion framework developed in the ProveIt coefficient-calculus monograph
\cite{ProveItCCC} and the Lambert-core inversion framework of the companion
series-and-transseries notes \cite{ProveItST,ProveItGammaBarnes}.

\begin{warningbox}
In this article, ``the inverse double factorial'' means a \emph{compositional}
large real inverse of a specified analytic interpolation.  The two parity
branches give two canonical inverses once parity is fixed.  The OEIS cosine
interpolation gives a third, eventually monotone inverse with a periodic
transseries sector.  The reciprocal \(1/(x!!)\) is different; its all-orders
coefficients are also recorded in \cref{sec:multiplicative-coefficients}.
\end{warningbox}

The word \emph{complete} refers to all orders of the positive-real
algebraic--logarithmic transseries, with general coefficient formulae and
fixed-order analytic remainders.  The Stirling series is divergent.  Its
terminant or resurgent completion is a separate hyperasymptotic layer; the
transport rule for such beyond-all-orders terms is given in
\cref{sec:beyond-all-orders}.

\section{Exact continuations and the parity issue}
\label{sec:continuations}

\subsection{The two exact parity branches}

Let \(\delta\in\{0,1\}\), where \(\delta=0\) denotes the even subsequence and
\(\delta=1\) the odd subsequence.  Define
\begin{equation}
 \kappa_\delta=\left(\frac2\pi\right)^{\delta/2},
 \qquad
 \DF_\delta(x)=\kappa_\delta 2^{x/2}\Gamma\!\left(\frac x2+1\right).
 \label{eq:parity-branches}
\end{equation}
For every nonnegative integer \(n\equiv\delta\pmod2\),
\begin{equation}
 \DF_\delta(n)=n!!.
 \label{eq:branch-interpolation}
\end{equation}
Indeed, \((2m)!!=2^m m!\), while
\((2m-1)!!=2^m\Gamma(m+1/2)/\sqrt\pi\).  Equivalently, for integer \(n\),
\begin{equation}
 n!!=2^{n/2}\Gamma\!\left(\frac n2+1\right)
 \left(\frac2\pi\right)^{(1-(-1)^n)/4}.
 \label{eq:integer-unified-exact}
\end{equation}
The normalization used repeatedly below is
\begin{equation}
 C_\delta=\sqrt\pi\,\kappa_\delta
 =\begin{cases}\sqrt\pi,&\delta=0,\\[0.2em]\sqrt2,&\delta=1.\end{cases}
 \label{eq:Cdelta}
\end{equation}

\subsection{The cosine interpolation recorded by OEIS}

A smooth interpolation that agrees with both parity classes is
\begin{equation}
 \Dc(x)=2^{x/2}\Gamma\!\left(\frac x2+1\right)
 \left(\frac2\pi\right)^{(1-\cos\pi x)/4}.
 \label{eq:cosine-interpolation}
\end{equation}
This is equivalent to the OEIS formula with
\(h=x/2\) and exponent \(\sin^2(\pi h)/2\) \cite{OEIS}.  It is only one possible
interpolation, but it is particularly natural because the parity selector is
an entire periodic function.

The logarithmic derivative is
\begin{equation}
 \frac{\Dc'(x)}{\Dc(x)}
 =\frac12\log2+\frac12\psi\!\left(\frac x2+1\right)
 +\frac\pi4\log\!\left(\frac2\pi\right)\sin\pi x.
 \label{eq:cosine-log-derivative}
\end{equation}
Since \(\psi(x/2+1)=\log(x/2)+O(x^{-1})\), this derivative is positive for all
sufficiently large \(x\).  Hence \(\Dc\) has a unique large real inverse,
although that inverse differs asymptotically from both parity-resolved
inverses.

\subsection{Why the centre \texorpdfstring{\(X=x+1\)}{X=x+1} is optimal}
\label{sec:optimal-centre}

Put
\begin{equation}
 X=x+1,
 \qquad t=\frac X2.
 \label{eq:center-variables}
\end{equation}
Then
\begin{equation}
 \DF_\delta(x)=\kappa_\delta 2^{t-1/2}\Gamma(t+1/2).
 \label{eq:centered-gamma}
\end{equation}
The half-shift in \(\Gamma(t+1/2)\) annihilates every even inverse power in its
logarithmic Stirling series.  This sparse support is the reason the inverse
contains only odd inverse powers of its Lambert core.

\section{Coefficient toolkit}
\label{sec:toolkit}

\subsection{Bernoulli numbers and the half-shift identity}

The Bernoulli polynomials are defined by
\begin{equation}
 \frac{u\e^{zu}}{\e^u-1}
 =\sum_{m\ge0}B_m(z)\frac{u^m}{m!},
 \qquad B_m=B_m(0).
 \label{eq:bernoulli-generating}
\end{equation}
The value at the half-shift is
\begin{equation}
 B_m\!\left(\frac12\right)=(2^{1-m}-1)B_m.
 \label{eq:bernoulli-half}
\end{equation}
For odd \(m>1\), both sides vanish.

\subsection{Exponential Bell polynomials}

For \(0\leq m\leq n\), the partial exponential Bell polynomial is
\begin{equation}
 \BellP_{n,m}(x_1,x_2,\ldots)
 =n!\!\sum_{\substack{j_1+j_2+\cdots=m\\
                      j_1+2j_2+\cdots=n}}
 \prod_{r\ge1}\frac1{j_r!}\left(\frac{x_r}{r!}\right)^{j_r}.
 \label{eq:partial-bell}
\end{equation}
The complete polynomial is
\begin{equation}
 \BellC_n(x_1,\ldots,x_n)=\sum_{m=0}^n\BellP_{n,m}(x_1,\ldots),
 \qquad \BellC_0=1.
 \label{eq:complete-bell}
\end{equation}
The two generating identities needed below are
\begin{align}
 \exp\!\left(\sum_{r\ge1}x_r\frac{z^r}{r!}\right)
 &=\sum_{n\ge0}\BellC_n(x_1,\ldots,x_n)\frac{z^n}{n!},
 \label{eq:complete-bell-egf}\\
 \left(\sum_{r\ge1}u_rz^r\right)^m
 &=\sum_{n\ge m}\frac{m!}{n!}
 \BellP_{n,m}(1!u_1,2!u_2,\ldots)z^n.
 \label{eq:partial-bell-power}
\end{align}
These formulae make every coefficient below a finite sum indexed by integer
partitions of the target order.

\subsection{Unsigned Stirling numbers of the first kind}

Write
\(\stirlingone{n}{k}\) for the unsigned Stirling number of the first kind,
defined by
\begin{equation}
 u(u+1)\cdots(u+n-1)=\sum_{k=0}^n\stirlingone{n}{k}u^k.
 \label{eq:stirling-first}
\end{equation}
They give the general polynomials in the large-argument expansion of Lambert
\(W\); see \cref{sec:outer-logs}.

\section{The complete forward transseries}
\label{sec:forward}

\subsection{Centered logarithmic form}

Define the dominant phase
\begin{equation}
 \Phi(X)=\frac X2(\log X-1).
 \label{eq:Phi}
\end{equation}

\begin{theorem}[Centered double-factorial Stirling expansion]
\label{thm:centered-forward}
For either \(\delta\in\{0,1\}\), as \(x\to+\infty\) with \(X=x+1\),
\begin{equation}
 \log\DF_\delta(x)
 \sim \log C_\delta+\Phi(X)+\sum_{k\ge1}\frac{c_k}{X^{2k-1}},
 \label{eq:forward-log}
\end{equation}
where
\begin{equation}
 c_k=\frac{2^{2k-1}B_{2k}(1/2)}{2k(2k-1)}
 =\frac{(1-2^{2k-1})B_{2k}}{2k(2k-1)}.
 \label{eq:ck}
\end{equation}
For every fixed \(N\ge0\), truncation after \(c_NX^{-(2N-1)}\) has remainder
\begin{equation}
 \log\DF_\delta(x)-\log C_\delta-\Phi(X)
 -\sum_{k=1}^{N}\frac{c_k}{X^{2k-1}}
 =O\!\left(X^{-(2N+1)}\right).
 \label{eq:forward-log-remainder}
\end{equation}
Here and below an empty sum is understood when \(N=0\).
\end{theorem}

\begin{proof}
The duplication formula gives
\begin{equation}
 \log\Gamma\!\left(t+\frac12\right)
 =\log\Gamma(2t)-\log\Gamma(t)+(1-2t)\log2+\frac12\log\pi.
 \label{eq:duplication-log}
\end{equation}
Apply the ordinary positive-real Stirling expansion, with its fixed-order
remainder, to the two gamma terms.  Their Bernoulli corrections differ by
\[
 \frac{B_{2k}}{2k(2k-1)}
 \bigl((2t)^{-(2k-1)}-t^{-(2k-1)}\bigr)
 =\frac{B_{2k}(1/2)}{2k(2k-1)t^{2k-1}}.
\]
Using \(t=X/2\) and multiplying by the factor
\(\kappa_\delta2^{t-1/2}\) in \cref{eq:centered-gamma} gives
\cref{eq:forward-log,eq:ck}.  The same subtraction of the two fixed-order
Stirling remainders proves \cref{eq:forward-log-remainder}; see also the
standard Gamma remainder theory in \cite{DLMF511,Olver,NemesGamma}.
\end{proof}

The first logarithmic coefficients are
\begin{equation}
 \begin{split}
 c_1&=-\frac1{12},\qquad
 c_2=\frac7{360},\qquad
 c_3=-\frac{31}{1260},\qquad
 c_4=\frac{127}{1680},\\
 c_5&=-\frac{511}{1188},\qquad
 c_6=\frac{1414477}{360360},\qquad
 c_7=-\frac{8191}{156}.
 \end{split}
 \label{eq:first-ck}
\end{equation}
Thus the compact all-orders form is
\begin{equation}
 \DF_\delta(x)
 \sim C_\delta\left(\frac X\e\right)^{X/2}
 \exp\!\left(\sum_{k\ge1}c_kX^{1-2k}\right).
 \label{eq:forward-exponential}
\end{equation}
All parity dependence has been compressed into \(C_\delta\).

\subsection{Unified form for the integer sequence}

For integer \(n\), put \(X=n+1\) and \(\chi_n=(-1)^n\).  Then
\begin{equation}
 C_n=\pi^{(1+\chi_n)/4}2^{(1-\chi_n)/4}
 =(2\pi)^{1/4}\left(\frac\pi2\right)^{\chi_n/4}.
 \label{eq:Cn}
\end{equation}
Consequently,
\begin{equation}
 n!!\sim C_n\left(\frac{n+1}{\e}\right)^{(n+1)/2}
 \exp\!\left(\sum_{k\ge1}\frac{c_k}{(n+1)^{2k-1}}\right).
 \label{eq:integer-forward-unified}
\end{equation}
This is a single parity-marked asymptotic formula for OEIS~A006882.

For comparison, expanding around \(x\) instead of \(x+1\) gives the denser
identity
\begin{equation}
 \log\DF_\delta(x)
 \sim \frac x2(\log x-1)+\frac12\log x+\log C_\delta
 +\sum_{k\ge1}
 \frac{2^{2k-1}B_{2k}}{2k(2k-1)x^{2k-1}}.
 \label{eq:uncentered-forward}
\end{equation}
Although the logarithmic tail still has odd powers, the extra
\(\tfrac12\log x\) complicates inversion.  The shift \(X=x+1\) absorbs it into
the dominant phase exactly.

\subsection{Multiplicative, reciprocal, and arbitrary-power coefficients}
\label{sec:multiplicative-coefficients}

Define the sparse sequence
\begin{equation}
 \lambda_r=
 \begin{cases}
 c_{(r+1)/2},&r\ \text{odd},\\
 0,&r\ \text{even}.
 \end{cases}
 \label{eq:lambda-r}
\end{equation}
For an arbitrary scalar \(\alpha\), set
\begin{equation}
 \exp\!\left(\alpha\sum_{r\ge1}\lambda_rX^{-r}\right)
 =\sum_{n\ge0}A_n^{(\alpha)}X^{-n}.
 \label{eq:Aalpha-def}
\end{equation}

\begin{proposition}[General multiplicative coefficients]
\label{prop:Aalpha}
For every \(n\ge0\),
\begin{align}
 A_n^{(\alpha)}
 &=\frac1{n!}\BellC_n\bigl(\alpha1!\lambda_1,
                \alpha2!\lambda_2,\ldots,\alpha n!\lambda_n\bigr)
 \label{eq:Aalpha-bell}\\
 &=\sum_{\substack{j_1+2j_2+\cdots=n}}
   \prod_{r\ge1}\frac{(\alpha\lambda_r)^{j_r}}{j_r!}.
 \label{eq:Aalpha-partition}
\end{align}
Equivalently,
\begin{equation}
 A_0^{(\alpha)}=1,
 \qquad
 nA_n^{(\alpha)}
 =\alpha\sum_{r=1}^n r\lambda_rA_{n-r}^{(\alpha)}.
 \label{eq:Aalpha-recurrence}
\end{equation}
\end{proposition}

\begin{proof}
\Cref{eq:Aalpha-bell} is the complete Bell generating identity
\cref{eq:complete-bell-egf}; expanding the Bell polynomial gives
\cref{eq:Aalpha-partition}.  Logarithmic differentiation of
\cref{eq:Aalpha-def} and coefficient comparison gives
\cref{eq:Aalpha-recurrence}.
\end{proof}

Taking \(\alpha=1\) gives
\begin{equation}
 \DF_\delta(x)
 \sim C_\delta\left(\frac X\e\right)^{X/2}
 \sum_{n\ge0}\frac{A_n^{(1)}}{X^n},
 \label{eq:forward-multiplicative}
\end{equation}
whereas \(\alpha=-1\) gives the reciprocal transseries
\begin{equation}
 \frac1{\DF_\delta(x)}
 \sim C_\delta^{-1}\left(\frac\e X\right)^{X/2}
 \sum_{n\ge0}\frac{A_n^{(-1)}}{X^n}.
 \label{eq:reciprocal-transseries}
\end{equation}
More generally, \(\alpha\) gives every fixed power \(\DF_\delta(x)^\alpha\)
after multiplying by
\(C_\delta^\alpha(X/\e)^{\alpha X/2}\).

The first forward coefficients are
\begin{equation}
 \begin{array}{c|rrrrrrr}
 n&0&1&2&3&4&5&6\\ \midrule
 A_n^{(1)}
 &1&-\frac1{12}&\frac1{288}&\frac{1003}{51840}
 &-\frac{4027}{2488320}&-\frac{5128423}{209018880}
 &\frac{168359651}{75246796800}
 \end{array}
 \label{eq:first-A}
\end{equation}
Because the logarithmic tail is odd under \(X^{-1}\mapsto-X^{-1}\), one in fact has the exact coefficient symmetry
\begin{equation}
 A_n^{(-\alpha)}=(-1)^nA_n^{(\alpha)}
 \qquad(n\ge0).
 \label{eq:Aalpha-sign-symmetry}
\end{equation}
In particular, the reciprocal coefficients are \(A_n^{(-1)}=(-1)^nA_n^{(1)}\) at every order.

\subsection{The cosine interpolation in forward form}
\label{sec:cosine-forward}

Let
\begin{equation}
 C_*=(2\pi)^{1/4},
 \qquad
 \rho=\frac12\log\!\left(\frac\pi2\right)>0.
 \label{eq:Cstar-rho}
\end{equation}
Since \(\cos\pi x=-\cos\pi X\), \cref{eq:cosine-interpolation} and
\cref{thm:centered-forward} give
\begin{equation}
 \Dc(x)
 \sim C_*\left(\frac X\e\right)^{X/2}
 \exp\!\left(-\frac\rho2\cos\pi X\right)
 \exp\!\left(\sum_{k\ge1}c_kX^{1-2k}\right).
 \label{eq:cosine-forward-transseries}
\end{equation}
The periodic factor is order one and must not be hidden inside an algebraic
remainder.  If desired, it has the exact Fourier--Bessel expansion
\begin{equation}
 \exp\!\left(-\frac\rho2\cos\theta\right)
 =I_0\!\left(\frac\rho2\right)
 +2\sum_{m\ge1}(-1)^mI_m\!\left(\frac\rho2\right)\cos(m\theta).
 \label{eq:bessel-periodic}
\end{equation}
At integral \(x=n\), \(\cos\pi X=-(-1)^n\), so the periodic normalization
reduces exactly to \(C_n/C_*\) in \cref{eq:Cn}.

\section{Parity-resolved functional inverses}
\label{sec:parity-inverse}

\subsection{Existence and exact Lambert core}

For each \(\delta\), the branch \(\DF_\delta(x)\) is positive and eventually
strictly increasing, because
\begin{equation}
 \frac{\DF_\delta'(x)}{\DF_\delta(x)}
 =\frac12\log2+\frac12\psi\!\left(\frac x2+1\right)
 =\frac12\log x+O(x^{-1}).
 \label{eq:branch-log-derivative}
\end{equation}
Let \(\DF_\delta^{-1}(y)\) denote its large real inverse.  The two continuations are exact scalar multiples,
\begin{equation}
 \DF_1(x)=\sqrt{\frac2\pi}\,\DF_0(x),
 \qquad
 \DF_1^{-1}(y)=\DF_0^{-1}\!\left(\sqrt{\frac\pi2}\,y\right),
 \label{eq:branch-inverse-scaling}
\end{equation}
which is the exact origin of the universal coefficient blocks below.

Put
\begin{equation}
 R_\delta=\log\!\left(\frac y{C_\delta}\right),
 \qquad
 w_\delta=\W_0\!\left(\frac{2R_\delta}{\e}\right),
 \label{eq:R-w}
\end{equation}
\begin{equation}
 T_\delta=\frac{2R_\delta}{w_\delta}=\e^{w_\delta+1},
 \qquad
 \Lambda_\delta=\log T_\delta=w_\delta+1.
 \label{eq:T-Lambda}
\end{equation}
Then
\begin{equation}
 \Phi(T_\delta)=R_\delta,
 \qquad
 T_\delta(\Lambda_\delta-1)=2R_\delta,
 \qquad
 \Phi'(T_\delta)=\frac{\Lambda_\delta}{2}.
 \label{eq:core-identities}
\end{equation}
The principal branch \(W_0\) is the only real Lambert branch needed because
\(2R_\delta/\e>0\) for large \(y\).

For notational economy, the subscripts \(\delta\) on \(R,w,T,\Lambda\) are
suppressed in coefficient calculations below.

\subsection{Normalized displacement equation}

Write
\begin{equation}
 X=T(1+\epsilon),
 \qquad q=T^{-1},
 \qquad z=q^2.
 \label{eq:epsilon-def}
\end{equation}
Subtracting the core equation \(\Phi(T)=R\) from the logarithmic forward phase
and multiplying by \(2/T\) gives the exact formal normal form
\begin{equation}
 \Lambda\epsilon+(1+\epsilon)\log(1+\epsilon)-\epsilon
 +2\sum_{k\ge1}c_kz^k(1+\epsilon)^{-(2k-1)}=0.
 \label{eq:normalized-inverse}
\end{equation}
Because the perturbation uses only powers \(z^k=q^{2k}\), its unique formal
solution has the form
\begin{equation}
 \epsilon=E(z)=\sum_{n\ge1}p_n(\Lambda)z^n.
 \label{eq:E-series}
\end{equation}
Since \(X-T=T\epsilon\), this produces odd powers
\(T^{-(2n-1)}\) in the absolute displacement.

\begin{mainbox}
For \(\delta\in\{0,1\}\), define \(R_\delta,T_\delta,\Lambda_\delta\) by
\cref{eq:R-w,eq:T-Lambda}.  Then for every fixed \(N\ge0\),
\begin{equation}
 \boxed{
 \DF_\delta^{-1}(y)
 =-1+T_\delta+
 \sum_{n=1}^{N}\frac{p_n(\Lambda_\delta)}{T_\delta^{2n-1}}
 +O\!\left(\frac1{T_\delta^{2N+1}\Lambda_\delta}\right)}
 \label{eq:inverse-main}
\end{equation}
as \(y\to+\infty\).  All parity dependence is contained in
\(R_\delta=\log(y/C_\delta)\); the coefficient functions \(p_n\) are universal.
\end{mainbox}

\subsection{Triangular coefficient recurrence}
\label{sec:inverse-recurrence}

Let
\begin{equation}
 E_{n-1}(z)=\sum_{j=1}^{n-1}p_j(\Lambda)z^j.
 \label{eq:Etrunc}
\end{equation}

\begin{theorem}[All-orders inverse coefficient engine]
\label{thm:inverse-recurrence}
The coefficients in \cref{eq:inverse-main} are uniquely determined by
\begin{equation}
 \boxed{
 p_n(\Lambda)=-\frac1\Lambda\coeff{z}{n}
 \left\{
 (1+E_{n-1})\log(1+E_{n-1})-E_{n-1}
 +2\sum_{k=1}^{n}c_kz^k(1+E_{n-1})^{-(2k-1)}
 \right\}.}
 \label{eq:pn-coeff-recurrence}
\end{equation}
Moreover,
\begin{equation}
 p_n(\Lambda)\in\Lambda^{-1}\Q[\Lambda^{-1}],
 \qquad
 \deg_{\Lambda^{-1}}p_n\le2n-1,
 \label{eq:pn-ring}
\end{equation}
and its two extreme coefficients are
\begin{equation}
 p_n(\Lambda)
 =-\frac{2c_n}{\Lambda}+O(\Lambda^{-2}),
 \qquad
 [\Lambda^{-(2n-1)}]p_n
 =\frac1{3^n}\binom{1/2}{n}.
 \label{eq:pn-extremes}
\end{equation}
\end{theorem}

\begin{proof}
At order \(z^n\), the only occurrence of the unknown \(p_n\) in
\cref{eq:normalized-inverse} is \(\Lambda p_nz^n\).  All other terms involve
only \(p_1,\ldots,p_{n-1}\), which proves existence, uniqueness, and
\cref{eq:pn-coeff-recurrence}.

The nonlinear kernel starts at \(E^2/2\), and all binomial coefficients in the
Bernoulli sector are rational.  Inductively, the coefficient inside braces is
a polynomial in \(\Lambda^{-1}\) of degree at most \(2n-2\); division by
\(\Lambda\) proves \cref{eq:pn-ring}.  Its part independent of lower
coefficients is \(2c_n\), proving the first statement in
\cref{eq:pn-extremes}.

For the highest possible pole \(\Lambda^{-(2n-1)}\), only the quadratic kernel
and the first Bernoulli coefficient can contribute.  Since \(2c_1=-1/6\), the
maximal-pole generating equation is
\[
 \Lambda E+\frac12E^2-\frac16z=0.
\]
The solution with zero constant term is
\[
 E=-\Lambda+\sqrt{\Lambda^2+z/3}
 =\sum_{n\ge1}\binom{1/2}{n}
   \frac{z^n}{3^n\Lambda^{2n-1}},
\]
which proves the second statement.
\end{proof}

\subsection{Explicit Bell-polynomial formula}

The recurrence can be expanded into a finite combinatorial sum.  Use
\(\BellP_{0,0}=1\) and \(\BellP_{r,0}=0\) for \(r>0\).

\begin{corollary}[Partial-Bell form]
\label{cor:pn-bell}
For \(n\ge1\),
\begin{align}
 p_n(\Lambda)=-\frac1\Lambda\Bigg[{}
 &\sum_{m=2}^{n}\frac{(-1)^m(m-2)!}{n!}
 \BellP_{n,m}(1!p_1,2!p_2,\ldots)
 \notag\\
 &+2\sum_{k=1}^{n}c_k
   \sum_{m=0}^{n-k}
   \binom{1-2k}{m}\frac{m!}{(n-k)!}
   \BellP_{n-k,m}(1!p_1,2!p_2,\ldots)
 \Bigg].
 \label{eq:pn-bell}
\end{align}
Every \(p_j\) occurring on the right has \(j<n\).
\end{corollary}

\begin{proof}
Use
\begin{equation}
 (1+E)\log(1+E)-E
 =\sum_{m\ge2}\frac{(-1)^m}{m(m-1)}E^m,
 \label{eq:kernel-series}
\end{equation}
expand \((1+E)^{1-2k}\) binomially, and apply
\cref{eq:partial-bell-power}.  The factors
\(m!/[m(m-1)]=(m-2)!\) give the first line.
\end{proof}

\subsection{First five inverse blocks}
\label{sec:first-blocks}

The recurrence gives
\begin{align}
 p_1={}&\frac1{6\Lambda},
 \label{eq:p1}\\[0.2em]
 p_2={}&-\frac{14\Lambda^2+10\Lambda+5}{360\Lambda^3},
 \label{eq:p2}\\[0.2em]
 p_3={}&\frac{
 2232\Lambda^4+1176\Lambda^3+714\Lambda^2+350\Lambda+105}
 {45360\Lambda^5},
 \label{eq:p3}\\[0.2em]
 p_4={}&-\frac1{5443200\Lambda^7}
 \bigl(822960\Lambda^6+292536\Lambda^5+136956\Lambda^4
 \notag\\[-0.3em]
 &\hspace{7.2em}{}+68040\Lambda^3+32970\Lambda^2
 +12250\Lambda+2625\bigr),
 \label{eq:p4}\\[0.2em]
 p_5={}&\frac1{359251200\Lambda^9}
 \bigl(309052800\Lambda^8+77920128\Lambda^7+26024592\Lambda^6
 \notag\\[-0.3em]
 &\hspace{3.4em}{}+10019592\Lambda^5+4516028\Lambda^4
 +2023560\Lambda^3+812350\Lambda^2
 \notag\\[-0.3em]
 &\hspace{3.4em}{}+242550\Lambda+40425\bigr).
 \label{eq:p5}
\end{align}
Thus, through three correction blocks,
\begin{align}
 \DF_\delta^{-1}(y)={}&-1+T
 +\frac1{6T\Lambda}
 -\frac{14\Lambda^2+10\Lambda+5}{360T^3\Lambda^3}
 \notag\\
 &+\frac{2232\Lambda^4+1176\Lambda^3+714\Lambda^2
 +350\Lambda+105}{45360T^5\Lambda^5}
 +O\!\left(\frac1{T^7\Lambda}\right).
 \label{eq:inverse-three-blocks}
\end{align}

\subsection{Scaling from the half-shifted inverse Gamma problem}
\label{sec:gamma-scaling}

Let
\begin{equation}
 a_k=\frac{B_{2k}(1/2)}{2k(2k-1)}.
 \label{eq:ak-gamma}
\end{equation}
For the half-shifted inverse Gamma problem, the normalized equation has the same
kernel as \cref{eq:normalized-inverse}, with perturbation coefficients \(a_k\)
instead of \(2c_k\).  By \cref{eq:ck},
\begin{equation}
 2c_k=4^k a_k.
 \label{eq:scaling-perturbation}
\end{equation}
The coefficient recurrence is homogeneous with respect to the weight
\(\operatorname{wt}(a_k)=k\).  Therefore, if \(d_{2n-1}(\Lambda)\) denotes the
corresponding inverse Gamma block in the normalization of
\cite{ProveItGammaBarnes}, then
\begin{equation}
 \boxed{p_n(\Lambda)=4^n d_{2n-1}(\Lambda).}
 \label{eq:gamma-scaling}
\end{equation}
This identity both explains the displayed coefficients and supplies a strong
independent regression test for symbolic implementations.

\subsection{Analytic remainder proof}
\label{sec:inverse-remainder}

\begin{proof}[Proof of the remainder in \cref{eq:inverse-main}]
Substitute
\(
 E_N(z)=\sum_{n=1}^Np_n(\Lambda)z^n
\)
into \cref{eq:normalized-inverse}.  By construction, its formal residual is
\begin{equation}
 \Lambda E_N+(1+E_N)\log(1+E_N)-E_N
 +2\sum_{k=1}^{N+1}c_kz^k(1+E_N)^{-(2k-1)}
 =O(z^{N+1}).
 \label{eq:formal-residual}
\end{equation}
Restoring the factor \(T/2\) removed in the normalization turns this into a
forward logarithmic residual \(O(T^{-(2N+1)})\).  The omitted forward Stirling
terms have the same or smaller order by
\cref{eq:forward-log-remainder}.  Finally,
\begin{equation}
 \frac{\dd}{\dd X}\left(
 \Phi(X)+\sum_{k\ge1}c_kX^{1-2k}\right)
 =\frac12\log X+O(X^{-2})
 =\frac\Lambda2+o(1)
 \label{eq:inverse-slope}
\end{equation}
near \(X=T\).  The mean-value theorem divides the phase residual by a quantity
asymptotic to \(\Lambda/2\), yielding
\(O(T^{-(2N+1)}/\Lambda)\) for the inverse error.
\end{proof}

\section{The full outer iterated-logarithmic layer}
\label{sec:outer-logs}

Keeping \(W_0\) unexpanded is numerically and symbolically superior.  To embed
the result in a conventional polynomial-logarithmic transseries, however, set
\begin{equation}
 Z=\frac{2R}{\e},
 \qquad L_1=\log Z,
 \qquad L_2=\log L_1.
 \label{eq:L1L2}
\end{equation}
For \(Z\to+\infty\),
\begin{equation}
 W_0(Z)\sim L_1-L_2+\sum_{r\ge1}\frac{P_r(L_2)}{L_1^r},
 \label{eq:W-expansion}
\end{equation}
where the general Lambert polynomial is
\begin{equation}
 \boxed{
 P_r(t)=\sum_{m=1}^{r}
 (-1)^{r-m}\stirlingone{r}{r-m+1}\frac{t^m}{m!}.}
 \label{eq:Lambert-polynomial}
\end{equation}
The first three are
\begin{equation}
 P_1=t,
 \qquad P_2=\frac{t^2-2t}{2},
 \qquad P_3=\frac{2t^3-9t^2+6t}{6}.
 \label{eq:first-Lambert-polys}
\end{equation}
This is the standard all-orders expansion of Lambert \(W\)
\cite{CorlessEtAl,DLMF413}.

To expand \(T=2R/W_0(Z)\), define
\begin{equation}
 u_1=-L_2,
 \qquad u_j=P_{j-1}(L_2)\quad(j\ge2),
 \label{eq:uj}
\end{equation}
so that
\(
 W_0(Z)=L_1(1+\sum_{j\ge1}u_jL_1^{-j})
\).
Define reciprocal polynomials \(h_n(L_2)\) by
\begin{equation}
 h_0=1,
 \qquad
 h_n=-\sum_{j=1}^nu_jh_{n-j}.
 \label{eq:hn-recurrence}
\end{equation}
Then
\begin{equation}
 \boxed{
 T\sim\frac{2R}{L_1}\sum_{n\ge0}\frac{h_n(L_2)}{L_1^n},
 \qquad
 \Lambda=1+L_1-L_2+\sum_{r\ge1}\frac{P_r(L_2)}{L_1^r}.}
 \label{eq:T-outer-full}
\end{equation}
For example,
\begin{equation}
 T\sim\frac{2R}{L_1}\left[
 1+\frac{L_2}{L_1}
 +\frac{L_2^2-L_2}{L_1^2}
 +\frac{L_2^3-\tfrac52L_2^2+L_2}{L_1^3}
 +\cdots\right].
 \label{eq:T-first-outer}
\end{equation}

For completely mechanical outer composition, let
\(A(t)=1+\sum_{k\ge1}a_kt^k\) and
\(A(t)^\alpha=\sum_{n\ge0}b_n^{(\alpha)}t^n\).  Logarithmic
differentiation gives the universal power recurrence
\begin{equation}
 b_0^{(\alpha)}=1,
 \qquad
 n b_n^{(\alpha)}
 =\sum_{k=1}^{n}\bigl((\alpha+1)k-n\bigr)
 a_k b_{n-k}^{(\alpha)}.
 \label{eq:outer-power-recurrence}
\end{equation}
Apply it first to
\(W_0(Z)/L_1=1+\sum u_kL_1^{-k}\), and then to
\begin{equation}
 \frac{\Lambda}{L_1}
 =1+\frac{1-L_2}{L_1}
 +\sum_{k\ge2}\frac{P_{k-1}(L_2)}{L_1^k}.
 \label{eq:Lambda-normalized-outer}
\end{equation}
Finite convolution of the resulting power coefficients expands every monomial
\(T^{-(2m-1)}\Lambda^{-j}\).  Thus the inner blocks require no implicit
composition step after \cref{eq:outer-power-recurrence}.

\begin{proposition}[Nested transseries hierarchy]
\label{prop:hierarchy}
Substitution of \cref{eq:T-outer-full} into \cref{eq:inverse-main} gives a
unique formal expansion in the ordered hierarchy
\begin{equation}
 \frac{R}{L_1},\quad
 \frac{R\,P(L_2)}{L_1^2},\quad
 \frac{R\,P(L_2)}{L_1^3},\ldots
 \quad\succ\quad 1
 \quad\succ\quad
 T^{-1}\Lambda^{-j}
 \quad\succ\quad
 T^{-3}\Lambda^{-j}\quad\succ\cdots.
 \label{eq:hierarchy}
\end{equation}
At every prescribed monomial only finitely many Bell/Stirling partitions
contribute.  Therefore \cref{eq:Lambert-polynomial,eq:hn-recurrence,eq:pn-bell} constitute general coefficient formulae for the fully expanded
outer-and-inner transseries.
\end{proposition}

\begin{remark}
The constant shift \(-1\) is beyond every fixed order of the outer
\(L_1^{-1}\)-series because \(R/L_1^N\to\infty\) for every fixed \(N\).  The
first Bernoulli inverse block is smaller still:
\(
 p_1/(T\Lambda)\asymp(12R)^{-1}
\).
This ordering is obscured if all terms are written as a single informal series
in \(1/\log y\).
\end{remark}

\section{Inverse of the OEIS cosine interpolation}
\label{sec:cosine-inverse}

The periodic interpolation \(\Dc\) has a qualitatively new inverse sector.  Put
\begin{equation}
 R_*=\log\!\left(\frac y{C_*}\right),
 \qquad
 w_*=W_0\!\left(\frac{2R_*}{\e}\right),
 \qquad
 T=\frac{2R_*}{w_*},
 \qquad
 \Lambda=1+w_*.
 \label{eq:cosine-core}
\end{equation}
Set
\begin{equation}
 X=T+S,
 \qquad q=T^{-1},
 \qquad \theta=\pi T,
 \qquad
 G(u)=(1+u)\log(1+u)-u.
 \label{eq:cosine-S}
\end{equation}
The exact formal equation obtained from
\cref{eq:cosine-forward-transseries} is
\begin{align}
 \mathcal F(q,S)={}&
 \frac\Lambda2S+\frac{G(qS)}{2q}
 +\sum_{k\ge1}c_kq^{2k-1}(1+qS)^{-(2k-1)}
 \notag\\
 &-\frac\rho2\cos(\theta+\pi S)=0.
 \label{eq:cosine-F}
\end{align}
The quotient \(G(qS)/q\) is regular at \(q=0\), because \(G(u)=u^2/2+O(u^3)\).
During coefficient extraction in \(q\), the slowly varying \(\Lambda\) and the rapid phase
\(\theta=\pi T\) are treated as coefficient parameters; this is an oscillatory
logarithmic transseries rather than an ordinary power series with constant coefficients.

\subsection{The periodic carrier}

At \(q=0\), the displacement \(s_0\) is determined by
\begin{equation}
 s_0=\mu\cos(\theta+\pi s_0),
 \qquad
 \mu=\frac\rho\Lambda.
 \label{eq:carrier-equation}
\end{equation}
It is \(O(\Lambda^{-1})\), much larger than the first Bernoulli inverse block
\(O(T^{-1}\Lambda^{-1})\).

\begin{theorem}[All-orders periodic carrier]
\label{thm:periodic-carrier}
For sufficiently large \(\Lambda\), the solution of
\cref{eq:carrier-equation} tending to zero has the convergent local
Lagrange--B\"urmann expansion
\begin{equation}
 \boxed{
 s_0(\theta,\Lambda)
 =\sum_{m\ge1}\frac{\mu^m}{m!}
 \left.\frac{\dd^{m-1}}{\dd u^{m-1}}
 \cos^m(\theta+\pi u)\right|_{u=0}.}
 \label{eq:carrier-lagrange}
\end{equation}
Equivalently,
\begin{equation}
 s_0
 =\sum_{m\ge1}\frac{\mu^m}{m!2^m}
 \sum_{j=0}^{m}\binom mj
 \bigl(\ii\pi(m-2j)\bigr)^{m-1}
 \e^{\ii(m-2j)\theta}.
 \label{eq:carrier-fourier}
\end{equation}
The expression in \cref{eq:carrier-fourier} is real.  Its first terms are
\begin{equation}
 s_0=\mu\cos\theta
 -\frac\pi2\mu^2\sin2\theta
 -\frac{\pi^2}{8}\mu^3(\cos\theta+3\cos3\theta)
 +O(\mu^4).
 \label{eq:carrier-first}
\end{equation}
\end{theorem}

\begin{proof}
\Cref{eq:carrier-equation} has the form
\(s_0=\mu\phi(s_0)\) with \(\phi(u)=\cos(\theta+\pi u)\).  Lagrange inversion
gives \cref{eq:carrier-lagrange}.  Expanding
\[
 \cos^m z=2^{-m}\sum_{j=0}^m\binom mj\e^{\ii(m-2j)z}
\]
and differentiating gives \cref{eq:carrier-fourier}.  Pairing \(j\) with
\(m-j\) proves reality.  Direct evaluation for \(m=1,2,3\) gives
\cref{eq:carrier-first}.
\end{proof}

\subsection{All algebraic corrections around the carrier}

Seek
\begin{equation}
 S(q)=s_0+\sum_{r\ge1}s_r(\theta,\Lambda)q^r.
 \label{eq:S-series}
\end{equation}
The carrier Jacobian is
\begin{equation}
 J=\partial_S\mathcal F(0,s_0)
 =\frac12\left(\Lambda+\rho\pi\sin(\theta+\pi s_0)\right).
 \label{eq:J}
\end{equation}
It is nonzero for all sufficiently large \(\Lambda\).

\begin{theorem}[Periodic triangular recurrence]
\label{thm:periodic-recurrence}
Let
\(
 S_{r-1}=s_0+\sum_{j=1}^{r-1}s_jq^j
\).
Then every coefficient in \cref{eq:S-series} is uniquely determined by
\begin{equation}
 \boxed{
 s_r=-\frac1J\coeff{q}{r}\mathcal F(q,S_{r-1}),
 \qquad r\ge1.}
 \label{eq:sr-recurrence}
\end{equation}
In particular,
\begin{equation}
 s_1=\frac{1/12-s_0^2/4}{J}.
 \label{eq:s1}
\end{equation}
For fixed \(N\), the large inverse of the cosine interpolation is
\begin{equation}
 \boxed{
 \Dc^{-1}(y)
 =-1+T+s_0(\theta,\Lambda)
 +\sum_{r=1}^{N}\frac{s_r(\theta,\Lambda)}{T^r}
 +O\!\left(\frac1{T^{N+1}\Lambda}\right).}
 \label{eq:cosine-inverse-main}
\end{equation}
Each \(s_r\) can in turn be expanded to arbitrary order in
\(\mu=\rho/\Lambda\) by \cref{eq:carrier-lagrange} and ordinary Bell
composition.
\end{theorem}

\begin{proof}
At order \(q^r\), the missing term \(s_rq^r\) enters linearly with coefficient
\(J\); all remaining contributions depend only on earlier coefficients.  This
proves \cref{eq:sr-recurrence}.  At order \(q\),
\[
 \frac{G(qS)}{2q}=\frac q4s_0^2+O(q^2),
 \qquad
 \sum_{k\ge1}c_kq^{2k-1}(1+qS)^{-(2k-1)}=c_1q+O(q^2),
\]
so \(Js_1+s_0^2/4+c_1=0\); using \(c_1=-1/12\) gives
\cref{eq:s1}.  Truncation through \(q^N\) leaves a phase residual
\(O(q^{N+1})\).  Since the exact phase derivative is
\(\Lambda/2+O(1)\), the root error is
\(O(q^{N+1}/\Lambda)\), proving \cref{eq:cosine-inverse-main}.
\end{proof}

\begin{warningbox}
The periodic carrier in \cref{eq:carrier-lagrange} is not a property of the
discrete sequence alone.  It records the choice of the cosine interpolation
\cref{eq:cosine-interpolation}.  A different smooth interpolation agreeing at
all integers can change this sector while leaving every integer value and both
parity-resolved transseries unchanged.
\end{warningbox}

\section{Algorithms and exact symbolic computation}
\label{sec:algorithms}

\subsection{Forward coefficients}

\begin{algorithm}[Forward and reciprocal coefficient table]
\label{alg:forward}
Given a truncation order \(M\):
\begin{enumerate}
\item compute \(c_k\) from \cref{eq:ck} for
      \(1\le k\le\lceil M/2\rceil\);
\item set \(\lambda_{2k-1}=c_k\) and \(\lambda_{2k}=0\);
\item compute \(A_n^{(\alpha)}\) either from the complete Bell formula
      \cref{eq:Aalpha-bell} or the recurrence
      \cref{eq:Aalpha-recurrence}.
\end{enumerate}
Only exact rational arithmetic is required when \(\alpha\in\Q\).
\end{algorithm}

\subsection{Inverse blocks}

\begin{algorithm}[Parity-resolved inverse blocks]
\label{alg:inverse}
The following truncated-series procedure computes
\(p_1,\ldots,p_N\) exactly in \(\Q(\Lambda)\).
\begin{lstlisting}
E := 0
for n := 1,...,N:
    H := (1+E) log(1+E) - E
    for k := 1,...,n:
        H := H + 2 c_k z^k (1+E)^(-(2k-1))
    p_n := - coefficient(truncate_z(H,n), z^n) / Lambda
    E := E + p_n z^n
return p_1,...,p_N
\end{lstlisting}
At step \(n\), every product and composition may be truncated immediately
above \(z^n\).  The symbolic residual check is
\begin{equation}
 \Lambda E_N+(1+E_N)\log(1+E_N)-E_N
 +2\sum_{k=1}^{N}c_kz^k(1+E_N)^{-(2k-1)}
 =O(z^{N+1}).
 \label{eq:symbolic-residual-check}
\end{equation}
\end{algorithm}

A compact Wolfram Language implementation is:
\begin{lstlisting}[language=Mathematica]
Clear[c, p, z, lam, E];
c[k_Integer?Positive] :=
  (1 - 2^(2 k - 1)) BernoulliB[2 k]/(2 k (2 k - 1));

inverseBlocks[N_Integer?NonNegative] := Module[{E = 0, H, out = {}},
  Do[
    H = (1 + E) Log[1 + E] - E
      + Sum[2 c[k] z^k (1 + E)^(-(2 k - 1)), {k, 1, n}];
    p[n] = -Coefficient[Normal@Series[H, {z, 0, n}], z, n]/lam;
    p[n] = Factor[p[n]];
    E += p[n] z^n;
    AppendTo[out, p[n]],
    {n, 1, N}];
  out
]
\end{lstlisting}

\subsection{Stable evaluation}

For a target supplied as \(Y=\log y\), never form \(y\) when it would overflow.
For branch \(\delta\), compute
\begin{equation}
 R=Y-\log C_\delta,
 \qquad w=W_0(2R/\e),
 \qquad T=2R/w,
 \qquad\Lambda=w+1,
 \label{eq:stable-evaluation}
\end{equation}
then evaluate \cref{eq:inverse-main}.  A final Newton correction in the
logarithmic phase is
\begin{equation}
 x_{\rm new}=x-
 \frac{\log\DF_\delta(x)-Y}
 {\tfrac12\log2+\tfrac12\psi(x/2+1)}.
 \label{eq:newton}
\end{equation}
This avoids overflow and is quadratically convergent once the asymptotic
approximation is in the local basin.

\section{Numerical verification}
\label{sec:numerics}

All values in this section were computed with 80-digit arithmetic from
\cref{eq:parity-branches}; no rounded forward values of \(y\) were used.

\subsection{Forward series}

Let \(F_M(x)\) denote \cref{eq:forward-multiplicative} truncated after
\(A_M^{(1)}X^{-M}\).  The table gives
\(|F_M(x)/\DF_\delta(x)-1|\), with \(\delta\equiv x\pmod2\).
\begin{table}[htbp]
\centering
\small
\begin{tabular}{rcrrrrr}
\toprule
\(x\)&parity&\(M=0\)&\(M=1\)&\(M=3\)&\(M=5\)&\(M=9\)\\
\midrule
10&even&\(7.590\,10^{-3}\)&\(4.330\,10^{-5}\)&\(2.599\,10^{-7}\)&\(4.974\,10^{-9}\)&\(1.388\,10^{-11}\)\\
11&odd &\(6.957\,10^{-3}\)&\(3.538\,10^{-5}\)&\(1.751\,10^{-7}\)&\(2.782\,10^{-9}\)&\(5.454\,10^{-12}\)\\
20&even&\(3.974\,10^{-3}\)&\(9.988\,10^{-6}\)&\(1.432\,10^{-8}\)&\(6.756\,10^{-11}\)&\(1.314\,10^{-14}\)\\
21&odd &\(3.793\,10^{-3}\)&\(9.014\,10^{-6}\)&\(1.166\,10^{-8}\)&\(4.972\,10^{-11}\)&\(7.956\,10^{-15}\)\\
50&even&\(1.635\,10^{-3}\)&\(1.483\,10^{-6}\)&\(3.106\,10^{-10}\)&\(2.113\,10^{-13}\)&\(9.594\,10^{-19}\)\\
51&odd &\(1.604\,10^{-3}\)&\(1.424\,10^{-6}\)&\(2.861\,10^{-10}\)&\(1.866\,10^{-13}\)&\(7.799\,10^{-19}\)\\
\bottomrule
\end{tabular}
\caption{Relative errors of the centered multiplicative expansion.}
\label{tab:forward-errors}
\end{table}

\subsection{Inverse series}

Let \(y=\DF_\delta(x)\) exactly and let \(I_N(y)\) denote
\cref{eq:inverse-main} through \(p_N/T^{2N-1}\); \(I_0=-1+T\).
The following table gives \(|I_N(y)-x|\).
\begin{table}[htbp]
\centering
\small
\begin{tabular}{rcrrrrrr}
\toprule
\(x\)&parity&\(N=0\)&\(N=1\)&\(N=2\)&\(N=3\)&\(N=5\)\\
\midrule
10&even&\(6.307\,10^{-3}\)&\(1.644\,10^{-5}\)&\(1.610\,10^{-7}\)&\(3.689\,10^{-9}\)&\(1.139\,10^{-11}\)\\
11&odd &\(5.581\,10^{-3}\)&\(1.210\,10^{-5}\)&\(9.981\,10^{-8}\)&\(1.934\,10^{-9}\)&\(4.263\,10^{-12}\)\\
20&even&\(2.605\,10^{-3}\)&\(1.752\,10^{-6}\)&\(4.775\,10^{-9}\)&\(3.103\,10^{-11}\)&\(7.651\,10^{-15}\)\\
21&odd &\(2.450\,10^{-3}\)&\(1.496\,10^{-6}\)&\(3.717\,10^{-9}\)&\(2.204\,10^{-11}\)&\(4.523\,10^{-15}\)\\
50&even&\(8.311\,10^{-4}\)&\(8.979\,10^{-8}\)&\(4.194\,10^{-11}\)&\(4.716\,10^{-14}\)&\(3.443\,10^{-19}\)\\
51&odd &\(8.111\,10^{-4}\)&\(8.422\,10^{-8}\)&\(3.785\,10^{-11}\)&\(4.094\,10^{-14}\)&\(2.767\,10^{-19}\)\\
100&even&\(3.575\,10^{-4}\)&\(9.580\,10^{-9}\)&\(1.147\,10^{-12}\)&\(3.316\,10^{-16}\)&\(1.590\,10^{-22}\)\\
\bottomrule
\end{tabular}
\caption{Absolute errors of the parity-resolved inverse transseries.}
\label{tab:inverse-errors}
\end{table}

The near-geometric improvement at these moderate arguments is not a claim of
convergence.  At fixed \(x\), the Bernoulli-driven expansion eventually
diverges; optimal truncation or a Newton correction should be used in high
precision work.

\section{Beyond all algebraic orders}
\label{sec:beyond-all-orders}

The coefficients \(c_k\) inherit the factorial growth of Bernoulli numbers.
Since
\begin{equation}
 B_{2k}\sim(-1)^{k-1}\frac{2(2k)!}{(2\pi)^{2k}},
 \label{eq:bernoulli-growth}
\end{equation}
the least term of the centered series occurs for \(2k\) of order \(\pi X\),
and the optimally truncated positive-real remainder is exponentially small on
the scale \(\e^{-\pi X}\), up to algebraic factors.  Exponentially improved
Gamma expansions express this remainder by terminant functions
\cite{NemesGamma}.

The inverse transport is universal.  Suppose an exponentially improved
forward logarithmic phase is written
\begin{equation}
 \log\DF_\delta(x)
 =\log C_\delta+\Phi(X)+H_{\rm alg}(X)+\mathcal E(X),
 \label{eq:forward-exp-sector}
\end{equation}
where \(\mathcal E\) is beyond every algebraic power.  If
\(X_{\rm alg}\) is the algebraic inverse transseries, then the first
beyond-all-orders inverse correction is
\begin{equation}
 \Delta X
 =-\frac{\mathcal E(X_{\rm alg})}
 {\Phi'(X_{\rm alg})+H_{\rm alg}'(X_{\rm alg})}
 +O(\mathcal E^2)
 \sim-\frac{2\mathcal E(T)}{\Lambda}.
 \label{eq:inverse-exp-transport}
\end{equation}
Higher powers of \(\mathcal E\) follow from ordinary Lagrange--B\"urmann
reversion in the phase coordinate.  Thus an optimally truncated forward sector
of scale \(\e^{-\pi T}\) induces an inverse sector of scale
\(\e^{-\pi T}/\Lambda\).  This section supplies the bridge to a resurgent or
hyperasymptotic completion without pretending that the divergent Bernoulli
series is convergent.

\section{Logical status and interpretation}
\label{sec:status}

Three kinds of statement have been used.

\paragraph{Exact identities.}
Equations \cref{eq:parity-branches,eq:integer-unified-exact,eq:cosine-interpolation} are exact wherever the gamma factors are defined.
The Lambert identities \cref{eq:core-identities} are also exact.

\paragraph{Formal coefficient identities.}
The Bell formulae and triangular recurrences are identities in formal power
series over \(\Q(\Lambda)\), or over the ring of periodic functions of
\(\theta\) and formal \(\Lambda^{-1}\)-series in the cosine case.  Their
coefficient uniqueness is algebraic.

\paragraph{Poincar\'e asymptotics.}
The forward remainder follows from standard positive-real Stirling estimates.
Eventual monotonicity and the residual-to-root argument prove the stated
fixed-order inverse remainders.  No convergence of the infinite Bernoulli
series is asserted.

\section{Conclusions}
\label{sec:conclusions}

The double factorial has a particularly clean all-orders structure after the
single shift \(X=x+1\):
\begin{enumerate}
\item The even and odd branches share the same centered logarithmic tail;
      parity survives only in \(C_0=\sqrt\pi\) and \(C_1=\sqrt2\).
\item The forward coefficients are explicit Bernoulli combinations
      \(c_k=(1-2^{2k-1})B_{2k}/[2k(2k-1)]\), and every multiplicative,
      reciprocal, or fixed-power coefficient is a complete Bell polynomial.
\item The large inverse is controlled exactly by
      \(T=2R/W_0(2R/\e)\).  Its correction support is the odd grid
      \(T^{-1},T^{-3},\ldots\), and the universal blocks
      \(p_n(\Lambda)\) are generated by the finite recurrence
      \cref{eq:pn-coeff-recurrence} or the Bell formula
      \cref{eq:pn-bell}.
\item The complete outer logarithmic layer follows from the unsigned-Stirling
      formula \cref{eq:Lambert-polynomial} and the reciprocal recurrence
      \cref{eq:hn-recurrence}.
\item Merging parity classes by the OEIS cosine interpolation creates a new
      periodic carrier of order \(1/\Lambda\).  Its coefficients are given by
      the Lagrange/Fourier formula \cref{eq:carrier-fourier}; all subsequent
      powers of \(T^{-1}\) follow from \cref{eq:sr-recurrence}.
\end{enumerate}

\begin{summarybox}
For symbolic work, centre at \(X=x+1\), retain Lambert \(W_0\) exactly, generate
\(c_k\) from Bernoulli numbers, exponentiate with complete Bell polynomials,
and invert with the triangular \(z=T^{-2}\) recurrence.  Expand \(W_0\) into
iterated logarithms only as a final presentation layer.  When parity is not
specified, state the interpolation explicitly: its choice changes the inverse
transseries.
\end{summarybox}

\appendix

\section{A hand derivation of the first two inverse blocks}
\label{app:hand-blocks}

Since
\begin{equation}
 (1+E)\log(1+E)-E=\frac12E^2-\frac16E^3+\cdots,
 \label{eq:kernel-hand}
\end{equation}
the coefficient of \(z\) in \cref{eq:normalized-inverse} is
\[
 \Lambda p_1+2c_1=0,
\]
whence \(p_1=1/(6\Lambda)\).  At order \(z^2\),
\begin{equation}
 \Lambda p_2+\frac12p_1^2-2c_1p_1+2c_2=0.
 \label{eq:p2-hand-equation}
\end{equation}
Using \(c_1=-1/12\), \(c_2=7/360\), and \(p_1=1/(6\Lambda)\) gives
\[
 p_2=-\frac1\Lambda\left(
 \frac7{180}+\frac1{36\Lambda}+\frac1{72\Lambda^2}\right)
 =-\frac{14\Lambda^2+10\Lambda+5}{360\Lambda^3}.
\]
This illustrates why the recurrence produces a finite polynomial in
\(\Lambda^{-1}\) at every order.

\section{A sixth inverse block}
\label{app:p6}

For regression testing at higher order, the next coefficient is
\begin{align}
 p_6(\Lambda)=-\frac1{5884534656000\Lambda^{11}}
 \bigl({}&46195687238400\Lambda^{10}
 +8854370366400\Lambda^9
 +2171643318000\Lambda^8
 \notag\\
 &+605515022112\Lambda^7
 +212869408752\Lambda^6
 +84136852800\Lambda^5
 \notag\\
 &+35589153600\Lambda^4
 +14234019800\Lambda^3
 +4868113250\Lambda^2
 \notag\\
 &+1213962750\Lambda+165540375\bigr).
 \label{eq:p6}
\end{align}
Its leading coefficient is \(-2c_6/\Lambda\), and its highest pole equals
\(3^{-6}\binom{1/2}{6}\Lambda^{-11}\), as required by
\cref{eq:pn-extremes}.

\section{Notation dictionary}
\label{app:notation}

\begin{longtable}{>{\raggedright\arraybackslash}p{0.21\textwidth}
                  >{\raggedright\arraybackslash}p{0.70\textwidth}}
\toprule
Symbol & Meaning\\
\midrule
\endfirsthead
\toprule
Symbol & Meaning\\
\midrule
\endhead
\(\delta\) & Parity label: \(0\) for the even branch, \(1\) for the odd branch.\\
\(\DF_\delta\) & Exact gamma continuation of one parity subsequence.\\
\(\Dc\) & OEIS cosine interpolation agreeing with both parity classes.\\
\(X=x+1\) & Optimal centred variable.\\
\(C_0,C_1\) & Branch constants \(\sqrt\pi,\sqrt2\).\\
\(C_*\) & Mean cosine-interpolation constant \((2\pi)^{1/4}\).\\
\(\rho\) & Periodic amplitude parameter \(\tfrac12\log(\pi/2)\).\\
\(\Phi(X)\) & Dominant phase \(X(\log X-1)/2\).\\
\(c_k\) & Centered logarithmic Stirling coefficient in \cref{eq:ck}.\\
\(\lambda_r\) & Sparse embedding of \(c_k\) on odd indices.\\
\(A_n^{(\alpha)}\) & Multiplicative coefficient for the power \(\alpha\).\\
\(R\) & Logarithmic target after removal of the branch constant.\\
\(w\) & Lambert variable \(W_0(2R/\e)\).\\
\(T\) & Exact Lambert core \(2R/w=\e^{w+1}\).\\
\(\Lambda\) & Core slope parameter \(\log T=w+1\).\\
\(q,z\) & Inner grid variables \(q=T^{-1}\), \(z=q^2\).\\
\(p_n(\Lambda)\) & Universal parity-resolved inverse block multiplying
                     \(T^{-(2n-1)}\).\\
\(L_1,L_2\) & Outer logarithms \(\log(2R/\e)\), \(\log L_1\).\\
\(P_r\) & Lambert polynomial in \cref{eq:Lambert-polynomial}.\\
\(s_0\) & Periodic carrier for the cosine interpolation inverse.\\
\(s_r\) & Algebraic correction around the periodic carrier.\\
\bottomrule
\end{longtable}

\begin{thebibliography}{99}

\bibitem{OEIS}
OEIS Foundation Inc.,
\newblock The On-Line Encyclopedia of Integer Sequences, A006882,
\newblock \emph{Double factorials: \(n!!=1\cdot3\cdot5\cdots n\) or
\(2\cdot4\cdot6\cdots n\)},
\newblock \url{https://oeis.org/A006882}, accessed 3 September 2026.

\bibitem{Meserve}
B.~E. Meserve,
\newblock Double factorials,
\newblock \emph{American Mathematical Monthly} \textbf{55} (1948), 425--426.

\bibitem{CorlessEtAl}
R.~M. Corless, G.~H. Gonnet, D.~E.~G. Hare, D.~J. Jeffrey, and D.~E. Knuth,
\newblock On the Lambert \(W\) function,
\newblock \emph{Advances in Computational Mathematics} \textbf{5} (1996),
329--359,
\newblock \url{https://doi.org/10.1007/BF02124750}.

\bibitem{DLMF413}
NIST Digital Library of Mathematical Functions,
\newblock \S4.13, Lambert \(W\)-function,
\newblock \url{https://dlmf.nist.gov/4.13}, accessed 3 September 2026.

\bibitem{DLMF511}
NIST Digital Library of Mathematical Functions,
\newblock \S5.11, Asymptotic expansions of the Gamma function,
\newblock \url{https://dlmf.nist.gov/5.11}, accessed 3 September 2026.

\bibitem{NemesGamma}
G. Nemes,
\newblock Error bounds and exponential improvements for the asymptotic
expansions of the Gamma function and its reciprocal,
\newblock \emph{Proceedings of the Royal Society of Edinburgh Section A:
Mathematics} \textbf{145} (2015), 571--596,
\newblock \url{https://doi.org/10.1017/S0308210513001558}.

\bibitem{Olver}
F.~W.~J. Olver,
\newblock \emph{Asymptotics and Special Functions},
\newblock A K Peters, Wellesley, MA, 1997 reprint of the 1974 edition.

\bibitem{vanDerHoeven}
J. van der Hoeven,
\newblock \emph{Transseries and Real Differential Algebra},
\newblock Lecture Notes in Mathematics 1888, Springer, Berlin, 2006.

\bibitem{ProveItCCC}
V. Reshetnikov,
\newblock \emph{Combinatorial Coefficient Calculus: A Unified Formulae
Monograph}, ProveIt repository,
\newblock \url{https://github.com/VladimirReshetnikov/ProveIt/tree/main/Analysis/FabiusFunction/docs/semi-formalized-research-frontiers/drafts/combinatorial-coefficient-calculus/Combinatorial_Coefficient_Calculus},
accessed 3 September 2026.

\bibitem{ProveItST}
V. Reshetnikov,
\newblock \emph{Series and Transseries}, ProveIt research-frontier notes,
\newblock \url{https://github.com/VladimirReshetnikov/ProveIt/tree/main/Analysis/FabiusFunction/docs/semi-formalized-research-frontiers/drafts/series-and-transseries},
accessed 3 September 2026.

\bibitem{ProveItGammaBarnes}
V. Reshetnikov,
\newblock \emph{Asymptotic Inversion of the Gamma and Barnes \(G\)-Functions},
ProveIt repository, 2026,
\newblock \url{https://github.com/VladimirReshetnikov/ProveIt/tree/main/Analysis/FabiusFunction/docs/semi-formalized-research-frontiers/drafts/series-and-transseries/special-function-inversion/Asymptotic_Inversion_Gamma_Barnes_G}.

\end{thebibliography}

\end{document}
